\section{引言}
\label{sec:intro}

对运动表现的预测建模是科学中最古老的定量课题之一，体育统计文献汗牛充栋：从棋类与球类比赛的等级分系统，到对跑步与游泳个人纪录的预测。然而，大多数体育数据分析研究都建立在定制的、私有的或快速变化的数据集之上，导致结果难以横向比较。每篇论文往往自选数据划分、预处理与指标，因此两个“最先进”的数字常常并不在度量同一件事。\wcabench{} 在一个几乎被机器学习社区忽视的场景中解决这一问题：由世界魔方协会（\wca）记录的竞技速拧。

\subsection{动机}
\label{sec:motivation}

三个问题促使我们构建该基准。

\paragraph{研究碎片化。}已有针对 \wca{} 数据的工作分散在单一任务上。先前研究分别用核密度估计预测比赛名次、用回归模型拟合世界纪录轨迹、用极值理论估计人类极限 \citep{glickman_rating, tatem_sprint, evt_coles}。由于各自使用不同的划分、指标与预处理，其结论无法比较。

\paragraph{缺乏标准化评估。}现有的“CubeBench”类基准面向\emph{魔方求解}：它们评估语言模型或智能体能否规划出一串动作将打乱的魔方复原 \citep{cubebench, cube_bench_mllm}。这类基准很有价值，但与 \wca{} 的真实比赛记录无关，也无法回答“哪个模型能更准确地预测选手成绩”或“哪种方法能更可靠地估计 DNF 概率”等问题。据我们所知，目前尚无针对真实速拧数据的预测与推断标准化基准。

\paragraph{领域特定结构被忽视。}\wca{} 数据具有通用管线不处理的独特结构：跨越数年的选手纵向历史；17 个现役项目间的跨项目技能迁移；会改变平均成绩计算方式的轮次格式（\code{ao5}、\code{mo3}、best-of-3）；特殊状态码（$-1$ 表示 \DNF{}、$-2$ 表示 \DNS{}、$0$ 表示无成绩）；以及多盲项目的复合编码。例如，忽略“5 次去极值平均”规则的模型会系统性地误校准时间与名次预测。我们认为，这些约束理应在基准中获得一等公民式的处理，而非在每篇论文中临时拼凑。

\subsection{贡献}
\label{sec:contributions}

我们提出 \wcabench{}——一个面向 \wca{} 比赛数据的预测与推断基准。我们的贡献如下：

\begin{enumerate}
  \item \textbf{可复现的数据基础设施。} 基于列式（Parquet）的管线，解码全部领域特定数值，包括 \wca{} 多盲复合编码、含去极值的轮次平均，以及特殊状态码（第~\ref{sec:dataset} 节）。该管线可单命令复现，并附带说明数据来源、规模、划分、偏差与允许/禁止用途的数据卡。

  \item \textbf{五个形式化任务及其领域挑战。} 我们定义 \task{1} 成绩预测（回归）、\task{2} 名次预测（排序）、\task{3} DNF 预测（不平衡分类）、\task{4} 人类极限估计（极值外推）与 \task{5} 技能迁移分析（因果推断）。每个任务均以“输入/输出/指标”三元组、基线分类体系及其非平凡性来源的领域挑战加以刻画（第~\ref{sec:tasks} 节）。

  \item \textbf{防泄漏评估协议与分层报告框架。} 我们规定带冻结统计量的滚动窗口协议、四维分层方案（项目、水平、时间、地区）以及带效应量的统计检验工具箱，以避免对小幅提升的过度解读（第~\ref{sec:protocol} 节）。

  \item \textbf{21 个真实数字的参考基线。} 我们在官方导出上报告真实基线结果（第~\ref{sec:results} 节），覆盖六大方法家族——统计、树模型、深度、图、贝叶斯与因果——每个任务至少 4 个基线，并同时给出单次运行与三种子均值。例如，基于树模型的回归在 \task{1} 上取得 \maelog{} $=0.164$；\wca{} Psych Sheet、Plackett--Luce 模型与核密度模拟在 \task{2} 上并列取得 Kendall $\kendall=0.767$；基于树模型的分类在 \task{3} 上取得 \aucpr{} $=0.373$。我们刻意不宣称达到最先进水平：本基准的科学价值在于其标准化问题，而非某个新模型。
\end{enumerate}

\subsection{主要发现概览}
\label{sec:findings}

由于我们发布的是真实数字而非假设数字，本文也给出第一批实证观察，并同时报告单次运行与三种子均值（第~\ref{sec:results-multiseed} 节）。简单且设定良好的模型仍然出奇地有竞争力。在 \task{2} 上，Plackett--Luce 模型、核密度模拟与官方 Psych Sheet 在 seed-42 运行中给出的 Kendall $\kendall=0.767$ 精确到小数点后四位完全一致（三种子 $0.8132\pm0.0342$），而图神经网络（GNN）仅达到 $0.567$（$0.7440\pm0.1258$）。在 \task{1} 上，序列模型（LSTM，\maelog{} $=1.044$，三种子 $0.9436\pm0.0754$）远差于基于树模型的回归（\maelog{} $=0.164$，三种子 $0.1055\pm0.0417$）；在 \task{3} 上，最佳树模型（\aucpr{} $=0.373$）与历史 DNF 率先验（\aucpr{} $=0.359$）之间的差异\emph{并不}显著（$p=0.749$，Cohen's $d=0.002$），而无树模型的贝叶斯 Beta--Binomial 在多种子均值上反而最好（$0.2938\pm0.0570$）。这些结果共同表明，在此样本规模下，规则感知特征结合树模型优于通用深度与图架构。对于人类极限估计，指数衰减模型给出三阶极限 $3.74$ 秒，变点模型给出 $3.02$ 秒，最稳定的方法（高斯过程叠加极值理论）给出 $3.29$ 秒——均明显高于社区常引用的领域锚点 $2.37$ 秒，是外推风险的具体例证。最后，在技能迁移上，“可识别”项目对数量随识别假设趋严而单调下降：Spearman 相关标记 $413$ 对，因果森林 $312$ 对，工具变量 $270$ 对，而双重差分仅 $6$ 对（三种子均值 $2$），量化了朴素方法的高估程度。这些观察在第~\ref{sec:discussion} 节讨论。

\subsection{伦理与范围}
\label{sec:scope}

\wcabench{} 是评估基准，而非预测产品。它仅使用公开的“WCA 数据库导出”，绝不抓取私有数据。我们明确禁止赌博、博彩以及任何针对、排名或羞辱个体选手的用途，并提供退出机制（第~\ref{sec:limitations} 节）。范围之外的工作包括魔方求解/步数优化任务（由 CubeBench 等魔方求解基准覆盖 \citep{cubebench, cube_bench_mllm}）、LLM 空间推理评测，以及实时直播预测系统。
